%\setchapterpreamble[u]{\margintoc}
\chapter{Introduction}
%\labch{intro}


\section*{Motivation}
\addcontentsline{toc}{section}{Motivation}

Of all the macromolecules one might find in the cells of living organisms, proteins are the class which are overwhelmingly responsible for the execution of the cell's functions, operating within and beyond its borders. Proteins are at the heart of vital cell functions like metabolism, \textsc{dna} repair, cell polarity, and scaffolding.

To execute these myriad functions, proteins often participate in sometimes complex chains of chemical and physicochemical reactions with other proteins more broadly referred to as ``protein-protein interactions'' or \textsc{ppi}s. Understanding which proteins interact with other proteins is vital to our knowledge of \textit{how} cells carry-out their many functions as well as the ways in which they fail in the case of disease.

In fact, cancers and other diseases which are caused by protein mutations often result in a change in protein binding affinity, altering how those proteins interact with others \cite{Jemimah_Gromiha_2020}. Infection and diseases caused by viruses like \textsc{hiv} and \textsc{sars}-\textsc{c}o\textsc{v-2} are also carried out by interactions between host and viral proteins \cite{hiv_ppi, sars_cov2_ppi}. What's more, treatments for these diseases sometimes come in the form of so-called ``biologic'' therapies like small peptides, antibiotics, or m\textsc{rna} vaccines, which in turn carry out their functions through the same fundamental processes as \textsc{ppi}s do.

While the study of \textsc{ppi}s has immeasurably shaped our understanding of biology, the methods by which we've traditionally come to characterize them require immense investments of time, labour, and capital. These methods, which include yeast two-hybrid (\textsc{y2h}), affinity chromatography, and mass spectroscopy, are primarily conducted at the bench with the aid of specialized material, equipment, and expertise \cite{vidal_fields_2014, snider_2015, low_recent_2021}.

Consuming a fraction of the labour and capital resources, computational methods for identifying putative \textsc{ppi}s are a compelling methodology for complimenting the challenges posed by traditional experimental methods. A single \textsc{y2h} experiment, for instance, is conducted over multiple days \cite{y2h_protocol, mehla_mapping_2015}. Contrastingly, the interactions between a given protein and the entire human proteome can be conducted in approximately three minutes using the \textsc{intrepppid} method through the online \textsc{ppi}.bio portal \cite{intrepppid}. This is to make no mention of the discrepancies in capital investment between these two examples, with \textsc{y2h} requiring expensive specialist equipment and materials, while the \textsc{rapppid} \& \textsc{intrepppid} models runs on commodity computer hardware \cite{rapppid, intrepppid}.

Building upon decades of painstaking research, sizable databases of \textsc{ppi}s have been compiled across a growing number of organisms \cite{string_v12, intact}. These databases, in turn, have greatly enabled methods for identifying putative \textsc{ppi}s using data analysis and mining. Combined with the explosion of advances in the field of deep learning in recent years, it is increasingly viable for designers of \textsc{ppi} inference methods to leverage neural networks to increase benchmark accuracies.

Despite best efforts, accurate \textsc{ppi} inference that generalizes to all use-cases remains regrettably out of reach. Surprisingly prevalent and persistent challenges in the nature of training datasets and model architecture frustrate efforts to train \textsc{ppi} inference models which perform well on out-of-distribution examples \cite{park_marcotte, Hamp_Rost_2015, black_box}. While these hurdles have long been known to researchers, the challenges remain relevant and affect recent methods \cite{black_box}.

This thesis contains a collection of three studies aimed at improving \textsc{ppi} inference performance on out-of-distribution data, as well as characterizing new challenges on the horizon that threaten the generalization of tomorrow's \textsc{ppi} inference models when trained on datasets with maximum sequence identity thresholds ranging from 50\% to 90\%. 

\section*{Thesis Organization \& Contributions}
\addcontentsline{toc}{section}{Thesis Organization \& Contributions}

This thesis is organized into four additional chapters whose main technical contributions are:

\begin{itemize}
    \item \textbf{\sffamily Chapter 2} -- This thesis begins with a comprehensive review of the biological and computational literature which inform this work.
    \item \textbf{\sffamily Chapter 3} -- This chapter introduces \textsc{rapppid}, a simple model for \textsc{ppi} inference which relies on judicious applications of regularization to achieve state-of-the-art performance at the time of publication. This study also highlights that contemporaneous models fail to generalize to proteins outside the training set, and introduces tokenization typically applied to natural language to amino acids.
    \item \textbf{\sffamily Chapter 4} -- Expanding on the work done on \textsc{rapppid}, Chapter 4 reports on an investigation into leveraging orthology to further generalize \textsc{ppi} inference to unseen proteins belonging to evolutionarily-distant organisms. The resulting method, \textsc{intrepppid},  out-performs \textsc{rapppid} and other contemporaneous models when trained on human \textsc{ppi}s, and subsequently evaluated on other organisms. To do so, a novel ``orthologous locality'' loss term is defined to reshape the protein embedding latent space, ``pushing'' orthologous proteins close to one another and ``pulling'' non-orthologous proteins apart. 
    \item \textbf{\sffamily Chapter 5} -- In which we reveal a before-uncharacterized source of data leakage in \textsc{ppi} inference models which incorporate pre-trained protein large language models (p\textsc{llm}s). Furthermore, we demonstrate how careful attention to the creation of training datasets can result in p\textsc{llm}s which are both performant and do not leak data. Finally, failures to generalize across p\textsc{llm}-based and non-p\textsc{llm}-based are described.
\end{itemize}
\newpage

\section*{Publications \& Author Contributions}
\addcontentsline{toc}{section}{Publications \& Author Contributions}
\begin{itemize}
    \item \textbf{\sffamily Chapter 3}
    \begin{tcolorbox}
    Szymborski, Joseph, and Amin Emad. 2022. ``\textsc{rapppid}: Towards Generalizable Protein Interaction Prediction with \textsc{awd}-\textsc{lstm} Twin Networks.'' \textit{Bioinformatics} 38 (16): 3958–67.\\\textsc{doi}: \href{https://10.1093/bioinformatics/btac429}{10.1093/bioinformatics/btac429}.
    \end{tcolorbox}
    \textsc{a.e.} and \textsc{j.s.} conceived the study, designed the project and the algorithm, and wrote the article. \textsc{j.s.} implemented the pipeline and performed the statistical analyses of the results. All authors read and approved the final article.
    

    \item \textbf{\sffamily Chapter 4} 
    \begin{tcolorbox}
    Szymborski, Joseph, and Amin Emad. 2024. “\textsc{intrepppid}—an Orthologue-Informed Quintuplet Network for Cross-Species Prediction of Protein–Protein Interaction.” \textit{Briefings in Bioinformatics} 25 (5): bbae405. \textsc{doi}: \href{https://doi.org/10.1093/bib/bbae405}{10.1093/bib/bbae405}.
    \end{tcolorbox}

    \textsc{a.e.} and \textsc{j.s.} contributed to project conceptualization, methodology as well as writing, editing, and reviewing. \textsc{a.e.} contributed to supervision and funding acquisition. \textsc{j.s.} contributed to software, data curation, investigation and visualization. All authors read and approved the final article.
    
    \item \textbf{\sffamily Chapter 5} \\
    The manuscript in Chapter 5 is under-review and is currently available as a pre-print on the bioRxiv online pre-print repository.
    \begin{tcolorbox}
    Szymborski, Joseph, and Amin Emad. 2025. “A Flaw in Using Pre-Trained pLLMs in Protein-Protein Interaction Inference Models.” \textit{bioRxiv}. \textsc{doi}: \href{https://doi.org/10.1101/2025.04.21.649858}{10.1101/2025.04.21.649858}.
    \end{tcolorbox}

    \textsc{a.e.} and \textsc{j.s.} contributed to project conceptualization, methodology as well as writing, editing, and reviewing.
\textsc{a.e.} contributed to supervision and funding acquisition. \textsc{j.s.} contributed to software, data curation, investigation
and visualization. All authors read and approved the final article.
    

\end{itemize}
